\documentclass[11pt]{article}
\usepackage[margin=1in]{geometry}
\usepackage{amsmath}
\usepackage{xspace}
\usepackage{graphicx}
\usepackage{subfig}
\usepackage{hyperref}
\usepackage{lineno}
\usepackage{stackengine}
\usepackage{tabularx}
\usepackage{rotating}
\usepackage{braket}
\usepackage{placeins}
\usepackage{caption}
\usepackage{amssymb}
\usepackage{authblk}

\usepackage[colorinlistoftodos]{todonotes}
\renewcommand*{\Authand}{, }
%\linenumbers

\def \nue{\ensuremath{\nu_e}\xspace}
\def \nuebar{\ensuremath{\bar{\nu}_e}\xspace}
\def \numu{\ensuremath{\nu_\mu}\xspace}
\def \numubar{\ensuremath{\bar{\nu}_\mu}\xspace}

\newcommand\barparen[1]{\stackon[.1pt]{\ensuremath{#1}}{\brabar}}
\DeclareCaptionType{equ}[][]
\begin{document}
\title{Measurement of the Absolute Drell-Yan Cross Section in $p+p$ and $p+d$ interactions with a 120 GeV Proton Beam at Fermilab}
\author[1]{Chatura Kuruppu} 
%\author{Chatura Kuruppu} 

\affil[1]{New Mexico State University, Las Cruces, NM 88003}
%\affil{University of South Carolina, Columbia, South Carolina 29208, USA}
\date{\today}
\maketitle
\begin{abstract}
Drell-Yan (DY) dimuon production in $p+p$ and $p+d$ interactions provides a powerful, direct probe of the antiquark structure of the nucleon. While the $p+d/p+p$ cross-section ratio is famously sensitive to the $\bar{d}(x)/\bar{u}(x)$ flavor asymmetry, measurements of the absolute cross sections are essential for independently constraining the $\bar{d}(x)$ and $\bar{u}(x)$ parton distribution functions (PDFs) and for testing higher-order QCD predictions. The SeaQuest (E906) experiment at Fermilab utilized a 120 GeV proton beam incident on liquid hydrogen and deuterium targets to collect a large dataset of DY events, specializing in the high-Bjorken-$x$ kinematic region where antiquark distributions remain poorly constrained. This talk will report on the status of the ongoing analysis to extract the absolute DY cross sections for both reactions. Specifically, we will discuss the unfolding techniques applied to correct for detector effects and present preliminary unfolded results for the $p+p$ and $p+d$ doubly differential cross sections in terms of dimuon invariant mass and Feynman-x ($x_F$), as well as the singly differential cross section $d\sigma/dp_T$ in terms of the transverse momentum $p_T$ of the dimuon. These detailed kinematic measurements will furnish crucial new data for global PDF analyses and provide stringent tests of non-perturbative models of nucleon structure.

\vspace{1em}
\noindent \textit{This work was supported in part by US DOE grant DE-FG02-94ER40847.}
\end{abstract}
\end{document}